\documentclass{article}
\usepackage[utf8]{inputenc}
\usepackage{polski}
\usepackage[polish]{babel}
\usepackage{amsmath}
\usepackage{amssymb}
\usepackage{geometry}

\geometry{a4paper, margin=1in}

\title{Teoria obliczeń i złożoność obliczeniowa - Lista 13}
\author{Aleksandra Połacik}
\date{Październik 2025}

\begin{document}

\maketitle


\section{Zadanie 1}
$O|V|^3 * O (|E|)$ \\
$O(|V^3|*|E|)$

\section{Zadanie 2}

Każda determistyczna TM jest niedeterministyczna TM, a to oznacza że nasz język A Należy do P.
1. Niech A należy do P: \\
2. Z definicji A jest rozstrzygalny przez pewną deterministyczną TM M działającą w czasie wielomianowym. \\
3. Ponieważ każda deterministyczna TM jest NTM to ozn. że język A jest rozstrzygalny przez niedeterministyczna maszyna turinga M dzialajca w czasie wielomianowym. \\

\section{Zadanie 3}

a. M1 - funkcja rozstrzygająca jezyk A, t`n. w należy do A:
$M_1(w)*$ = "tak" gdyby w należy do A \\
$M_1(w)*$ = "nie" gdy w nie należy do A  \\

def M(w): \\
    if $M_1(w)$ = "tak": \\
        return "tak" \\
    if $M_2(w)$ = "tak":   \\
        return "tak" \\
    return "nie"


$M_1(w)$ działa w czasie $O(|w|^{k_1})$
$M_2(w)$ działa w czasie $O(|w|^{k_2})$ gdzie k1 i k2 są stałymi \\

Odpowiedź: $O\!\left(n^{\max(k_1, k_2)}\right)$ \\


b. Maszyna niedeterministycznie dzieli słowo wejściowe $w$ na dwie części $w = xy$. Następnie sprawdza (niedeterministycznie w czasie wielomianowym), czy $x \in A$ oraz czy $y \in B$. Jeśli oba warunki są spełnione, akceptuje.

a. M1 - funkcja rozstrzygająca jezyk A, t`n. w należy do A:
$M_1(w)*$ = "tak" gdyby w należy do A \\
$M_1(w)*$ = "nie" gdy w nie należy do A  \\

Czyli musimy podzielić słowo na x i y prefix ma należeć do A i sufix do B

np. abba , na 5 części mozna podzielic na n+1 sposobów zeby sprawdzić konkatenacje \\
n+1 należy do O(n)

Odpowiedź: $O\!\left(n*n^{\max(k_1, k_2)+1}\right)$ \\

\section{Zadanie 4}
Zamkniętość NP ze względu na operację gwiazdki ($A^*$)Dla słowa $w$, maszyna niedeterministycznie wybiera liczbę $k$ (liczbę fragmentów) oraz miejsca podziału słowa na $w = w_1 w_2 \dots w_k$17.

Następnie dla każdego fragmentu $w_i$ maszyna niedeterministycznie sprawdza, czy $w_i \in A$1818.Ponieważ podziałów jest skończenie wiele, a sprawdzenie każdego fragmentu zajmuje czas wielomianowy, cała procedura mieści się w niedeterministycznym czasie wielomianowym.



\end{document}